\documentclass[10pt,a4paper]{article}
\usepackage[utf8]{inputenc}
\usepackage[T1]{fontenc}
\usepackage[brazilian]{babel}
\usepackage[margin=1.7cm]{geometry}
\usepackage{booktabs}
\usepackage{enumitem}
\usepackage{xcolor}
\usepackage{tikz}
\usetikzlibrary{positioning,arrows.meta,shapes.multipart}
\usepackage[hidelinks]{hyperref}

\setlength{\parskip}{0.4em}
\setlength{\parindent}{0pt}
\setlist[itemize]{noitemsep, topsep=2pt, leftmargin=1.4em}
\setlist[enumerate]{itemsep=2pt, topsep=2pt, leftmargin=1.5em}
\sloppy

\definecolor{umlheader}{RGB}{230, 238, 248}
\definecolor{umlborder}{RGB}{60, 90, 130}

\begin{document}

% Cabeçalho Objetivo
\begin{center}
    {\Large \textbf{Relatório Técnico — Blackjack 21 (POO)}}\\[0.25em]
    {\small Universidade Federal de São Paulo (Unifesp) — Programação Orientada a Objetos}\\[0.4em]
    {\small
    \begin{tabular}{ll@{\hspace{2.2cm}}ll}
        \textbf{Breno Porto Pinheiro do Prado} & RA: 185196 &
        \textbf{Felipe Pereira Silva} & RA: 188204 \\
        \textbf{Guilherme Bento Ramos} & RA: 185226 & & \\
    \end{tabular}}
\end{center}
\vspace{-0.2em}
\hrule
\vspace{0.4em}

\section*{1. Visão Geral e Lógica do Sistema}
O projeto implementa o jogo de cartas \textbf{Blackjack (21)} em \textbf{C\# (.NET 10)} utilizando a biblioteca gráfica \textbf{Raylib-cs}. O objetivo do jogador é obter uma pontuação próxima ou igual a 21 sem ultrapassar esse limite, competindo contra a Banca (\textit{Dealer}).

\begin{itemize}
    \item \textbf{Valores das Cartas:} Cartas de 2 a 10 assumem seu valor nominal; figuras (\textit{Jack}, \textit{Queen}, \textit{King}) valem 10 pontos; o \textbf{Ás (\textit{Ace})} possui valor dinâmico: inicia valendo \textbf{11 pontos}, mas caso a soma da mão ultrapasse 21, seu valor é reduzido para \textbf{1 ponto} (subtraindo 10 da soma por Ás disponível).
    \item \textbf{Distribuição Inicial (\texttt{ReiniciarJogo}):} O baralho de 52 cartas é recriado e embaralhado pelo algoritmo de \textit{Fisher-Yates}. Jogador e Banca recebem 2 cartas cada, sendo a segunda carta da Banca mantida oculta (\texttt{IsFaceUp = false}) durante a vez do jogador.
    \item \textbf{Turno do Jogador (\texttt{PedirCarta}):} Adiciona uma nova carta à mão do jogador. Caso a pontuação exceda 21 pontos (\texttt{IsBust() == true}), ocorre estouro imediato: a carta oculta da banca é revelada e a vitória é atribuída à Banca.
    \item \textbf{Turno Automático e Decisão da Banca (\texttt{Parar} e \texttt{ShouldHit}):} Quando o jogador encerra sua jogada, a Banca revela sua carta oculta e toma decisões autônomas: como o empate favorece a Banca, ela compra cartas apenas enquanto sua pontuação for estritamente menor que a pontuação do jogador e menor que 21 (\texttt{score < playerScore \&\& score < 21}). Assim que empata ou supera o jogador, a Banca para imediatamente para não correr risco de estourar.
    \item \textbf{Critério de Vitória:} O Jogador vence caso a Banca ultrapasse 21 pontos. Caso a Banca não estoure, ela vence ao igualar (empate) ou superar a pontuação do jogador.
\end{itemize}

\section*{2. Diagrama de Classes UML}

\begin{center}
\begin{tikzpicture}[
    scale=0.92, transform shape,
    umlclass/.style={
        draw=umlborder,
        thick,
        rectangle split,
        rectangle split parts=3,
        rectangle split part fill={umlheader,white,white},
        font=\scriptsize\ttfamily,
        align=left,
        inner sep=3.5pt
    },
    umlenum/.style={
        draw=umlborder,
        thick,
        rectangle split,
        rectangle split parts=2,
        rectangle split part fill={umlheader,white},
        font=\scriptsize\ttfamily,
        align=left,
        inner sep=3.5pt
    },
    inherit/.style={-{Triangle[open,length=6pt,width=6pt]}, thick, draw=umlborder},
    comp/.style={{Diamond[length=7pt,width=5pt]}-, thick, draw=umlborder},
    assoc/.style={{Diamond[open,length=7pt,width=5pt]}-, thick, draw=umlborder},
    dep/.style={-{Stealth[length=5pt]}, dashed, thick, draw=umlborder}
]

% BlackjackGame (Topo Centro)
\node[umlclass] (game) at (0, 0) {
    \textbf{BlackjackGame}
    \nodepart{second}
    - \_deck : Deck\\
    - \_player : Player\\
    - \_dealer : Dealer\\
    - \_btnPedirCarta : Button\\
    - \_btnParar : Button\\
    - \_btnReiniciar : Button\\
    - \_isGameOver : bool
    \nodepart{third}
    + ReiniciarJogo() : void\\
    + PedirCarta() : void\\
    + Parar() : void\\
    + Update() : void\\
    + Draw(int, int) : void
};

% Button (Topo Direita)
\node[umlclass] (button) at (6.2, 0.3) {
    \textbf{Button}
    \nodepart{second}
    + Bounds : Rectangle\\
    + Text : string\\
    + Color : Color\\
    + IsEnabled : bool
    \nodepart{third}
    + IsClicked() : bool\\
    + Draw(Font) : void
};

% Deck (Meio Esquerda)
\node[umlclass] (deck) at (-5.6, -0.2) {
    \textbf{Deck}
    \nodepart{second}
    - \_cards : List<Card>\\
    - \_random : Random
    \nodepart{third}
    + Reset() : void\\
    + Shuffle() : void\\
    + Draw() : Card
};

% Player (Abaixo Centro)
\node[umlclass] (player) at (0, -3.9) {
    \textbf{Player}
    \nodepart{second}
    - \_cards : List<Card>\\
    + Name : string\\
    + Cards : IReadOnlyList<Card>
    \nodepart{third}
    + AddCard(Card) : void\\
    + ClearCards() : void\\
    + CalculateScore() : int\\
    + IsBust() : bool
};

% Dealer (Abaixo Direita)
\node[umlclass] (dealer) at (6.2, -3.9) {
    \textbf{Dealer}
    \nodepart{second}
    \nodepart{third}
    + Dealer()\\
    + RevealCards() : void\\
    + ShouldHit(int) : bool
};

% Card (Abaixo Esquerda)
\node[umlclass] (card) at (-5.6, -4.0) {
    \textbf{Card}
    \nodepart{second}
    + Suit : Suit\\
    + Rank : Rank\\
    + IsFaceUp : bool
    \nodepart{third}
    + GetValue() : int\\
    + operator ==, !=, <, >
};

% Enums (Rodapé)
\node[umlenum] (suit) at (-7.0, -6.9) {
    \textbf{<<enum>> Suit}
    \nodepart{second}
    Hearts, Clubs,\\
    Diamonds, Spades
};

\node[umlenum] (rank) at (-3.8, -6.9) {
    \textbf{<<enum>> Rank}
    \nodepart{second}
    Two .. Ten,\\
    Jack, Queen, King, Ace
};

% Relações
\draw[comp] (game.west) -- node[above, font=\scriptsize]{1} (deck.east);
\draw[comp] (game.east) -- node[above, font=\scriptsize]{3} (button.west);
\draw[comp] (game.south) -- node[left, font=\scriptsize]{1} (player.north);
\draw[comp] (game.south east) -- (dealer.north west);

\draw[comp] (deck.south) -- node[left, font=\scriptsize]{0..52} (card.north);
\draw[assoc] (player.west) -- node[above, font=\scriptsize]{0..*} (card.east);
\draw[inherit] (dealer.west) -- node[above, font=\scriptsize]{herança} (player.east);

\draw[dep] (card.south west) -- (suit.north);
\draw[dep] (card.south east) -- (rank.north);

\end{tikzpicture}
\end{center}

\section*{3. Explicação das Classes e Conceitos de POO Aplicados}

\begin{enumerate}
    \item \textbf{\texttt{Card} (\texttt{src/Models/Card.cs}):} Representa uma carta individual. Utiliza \textbf{Encapsulamento} com propriedades somente leitura (\texttt{Suit} e \texttt{Rank}) inicializadas via \textbf{Construtor}, mantendo mutável apenas o estado \texttt{IsFaceUp}. Aplica \textbf{Sobrecarga de Operadores} (\texttt{==}, \texttt{!=}, \texttt{<}, \texttt{>}, \texttt{<=}, \texttt{>=}) para permitir comparações diretas de pontuação entre objetos \texttt{Card}.
    \item \textbf{\texttt{Deck} (\texttt{src/Models/Deck.cs}):} Representa o baralho de 52 cartas. Encapsula uma lista privada \texttt{List<Card> \_cards}, expondo apenas operações de alto nível: \texttt{Reset()} (recria as 52 combinações de naipe e valor), \texttt{Shuffle()} (embaralhamento) e \texttt{Draw()} (retira uma carta e recria o baralho automaticamente caso fique vazio).
    \item \textbf{\texttt{Player} (\texttt{src/Models/Player.cs}):} Representa o jogador e gerencia sua mão de cartas. Aplica \textbf{Encapsulamento estrito} ao manter \texttt{\_cards} privado e expor apenas uma visão somente leitura (\texttt{IReadOnlyList<Card> Cards}), impedindo modificações externas fora de \texttt{AddCard()} e \texttt{ClearCards()}. Concentra a regra de negócio de soma e ajuste dinâmico de múltiplos Ases em \texttt{CalculateScore()}.
    \item \textbf{\texttt{Dealer} (\texttt{src/Models/Dealer.cs}):} Representa a Banca. Demonstra \textbf{Herança} (\texttt{Dealer : Player}) e reuso de código, herdando toda a estrutura de cartas e cálculo de pontuação de \texttt{Player}, acionando o construtor da superclasse via \texttt{base("Banca")} e especializando o comportamento com \texttt{RevealCards()} e \texttt{ShouldHit(int playerScore)}.
    \item \textbf{\texttt{Button} (\texttt{src/UI/Button.cs}):} Encapsula a geometria (\texttt{Rectangle Bounds}), estado de ativação (\texttt{IsEnabled}), detecção de clique via colisão de mouse (\texttt{IsClicked()}) e renderização gráfica dos botões da interface.
    \item \textbf{\texttt{BlackjackGame} (\texttt{src/Game/BlackjackGame.cs}):} Classe controladora da partida. Demonstra \textbf{Composição e Associação entre classes} ao coordenar instâncias de \texttt{Deck}, \texttt{Player}, \texttt{Dealer} e os três objetos \texttt{Button}, separando a atualização lógica (\texttt{Update()}) da renderização visual (\texttt{Draw()}).
\end{enumerate}

\section*{4. Exemplos de Uso e Divisão do Projeto}

Durante a execução da aplicação gráfica (1280$\times$720 a 60 FPS), o usuário interage com os três botões principais:
\begin{itemize}
    \item \textbf{Botão ``Pedir Carta'':} Compra uma carta do baralho e atualiza instantaneamente o placar do jogador. Caso o jogador tenha um Ás (11) e um 8 (total 19) e compre um 7, o sistema recalcula automaticamente o Ás para 1 ponto (total 16). Se ultrapassar 21, os botões de ação são desabilitados e a derrota por estouro é exibida.
    \item \textbf{Botão ``Parar'':} Encerra a vez do jogador, revela a segunda carta da Banca e executa o laço de compra inteligente da Banca até que ela alcance a pontuação do jogador ou ultrapasse 21, exibindo o vencedor no painel central.
    \item \textbf{Botão ``Reiniciar Jogo'':} Limpa a mesa, reembaralha o baralho e inicia uma nova rodada.
\end{itemize}

\vspace{0.3em}
\begin{table}[h!]
\centering
\small
\begin{tabular}{lll}
\toprule
\textbf{Integrante} & \textbf{Arquivos Desenvolvidos} & \textbf{Foco Principal} \\
\midrule
Breno Porto P. do Prado (185196) & \texttt{Card.cs}, \texttt{Deck.cs} & Modelo de carta, sobrecarga e baralho \\
Felipe Pereira Silva (188204) & \texttt{Player.cs}, \texttt{Dealer.cs} & Pontuação/Ás, herança e decisão da banca \\
Guilherme Bento Ramos (185226) & \texttt{BlackjackGame.cs}, \texttt{Button.cs} & Fluxo de turnos, regras finais e interface \\
\bottomrule
\end{tabular}
\end{table}

\end{document}
